% !TeX root = part6.tex
\documentclass[../final.tex]{subfiles}

\begin{document}

% Линии фейнмановских диаграмм: стрелка фермиона и волнистый фотон.
\tikzset{
    feyn fermion/.style={
        draw=rtext,
        line width=0.9pt,
        postaction={decorate},
        decoration={
            markings,
            mark=at position 0.55 with {\arrow{Latex}}
        }
    },
    feyn photon/.style={
        draw=rtext,
        line width=0.85pt,
        decorate,
        decoration={
            snake,
            amplitude=1.6pt,
            segment length=6pt
        }
    }
}

\practice{6}{08.10.2026}{Квантовая электродинамика}

\rtheory{Сечение процесса \texorpdfstring{$AB\to CD$}{AB → CD}}

Рассматриваем процессы
\begin{equation*}
    e^+e^-\longrightarrow\mu^+\mu^-,
    \qquad
    \pi^+\pi^-\longrightarrow\pi^+\pi^-.
\end{equation*}
\rassumption{Будем считать, что это элементарные точечные частицы},
и для начала \rassumption{пренебрегаем спином: $s=0$}.
% В записи далее идёт пункт «2) … = 0», первый символ не читается.

Дифференциальное сечение записывается как
\begin{equation*}
    d\sigma=\frac{|\mathcal M|^2}{F}\,dQ,
\end{equation*}
где $\mathcal M$ --- \rconcept{инвариантная амплитуда},
\begin{equation*}
    F=|\vec v_{\text{отн}}|\,2E_A\,2E_B
\end{equation*}
--- \rconcept{начальная плотность потока}, а
\begin{equation*}
    dQ=(2\pi)^4\delta^{(4)}(p_C+p_D-p_A-p_B)\,
    \frac{d^3p_C}{(2\pi)^32E_C}\,
    \frac{d^3p_D}{(2\pi)^32E_D}
\end{equation*}
--- \rconcept{элемент фазового объёма}. Здесь $\delta^{(4)}$ ---
четырёхмерная дельта-функция, $p_A,\dots,p_D$ --- 4-векторы.

В системе центра масс получается
\begin{equation*}
    \left.\frac{d\sigma}{d\Omega}\right|_{\text{с.ц.м.}}
    =\frac{1}{64\pi^2s}\,\frac{p_f}{p_i}\,|\mathcal M|^2,
    \qquad
    s=(E_A+E_B)^2,
\end{equation*}
\begin{equation*}
    |\vec p_A|=|\vec p_B|=p_i,
    \qquad
    |\vec p_C|=|\vec p_D|=p_f.
\end{equation*}

\task[1]{Сечение в системе центра масс}
\label{task:cm_cross_section_sem6}

\condition
Получить выражение для $d\sigma/d\Omega$ в системе центра масс.

\solution
\subsection*{Фазовый объём}

В системе центра масс $\vec p_A+\vec p_B=0$ и полная энергия равна $\sqrt s$,
поэтому
\begin{equation*}
    \delta^{(4)}=\delta\!\left(E_C+E_D-\sqrt s\right)\cdot
    \delta^{(3)}\!\left(\vec p_C+\vec p_D\right).
\end{equation*}
Интегрируем по $d^3p_C$ с помощью трёхмерной дельта-функции:
\begin{equation*}
    dQ=\frac{(2\pi)^{4-6}}{4}\,
    \delta\!\left(E_C+E_D-\sqrt s\right)
    \frac{d^3p_D}{E_CE_D}
    =\frac{1}{16\pi^2}\,
    \delta\!\left(E_C+E_D-\sqrt s\right)
    \frac{d^3p_D}{E_CE_D}.
\end{equation*}
В сферических координатах
\begin{equation*}
    d^3p_D=p_D^2\,dp_D\,\underbrace{\sin\theta\,d\theta\,d\varphi}_{d\Omega}.
\end{equation*}

\subsection*{Переход к дельта-функции от энергии}

Из $E^2=p^2+m^2$ и $|\vec p_C|=|\vec p_D|=p_f$:
\begin{equation*}
    d(E_C+E_D)
    =d\!\left(\sqrt{p_f^2+m_C^2}+\sqrt{p_f^2+m_D^2}\right)
    =p_f\,dp_f\,\frac{E_C+E_D}{E_CE_D}.
\end{equation*}
Тогда
\begin{equation*}
    dQ=\frac{p_D\,d\Omega}{16\pi^2}
    \int\delta\!\left(E_C+E_D-\sqrt s\right)
    \frac{d(E_C+E_D)}{E_C+E_D}
    =\frac{p_D\,d\Omega}{16\pi^2}\cdot\frac{1}{\sqrt s}.
\end{equation*}

\subsection*{Поток}

Относительная скорость встречных частиц
$|\vec v_{\text{отн}}|=p_i/E_A+p_i/E_B$, поэтому
\begin{equation*}
    F=|\vec v_{\text{отн}}|\,2E_A\cdot2E_B
    =4p_i(E_A+E_B)
    =4p_i\sqrt s.
\end{equation*}

\answer
\begin{equation*}
    \frac{d\sigma}{d\Omega}
    =\frac{|\mathcal M|^2}{4p_i\sqrt s}\cdot
    \frac{p_f}{16\pi^2\sqrt s}
    =\frac{1}{64\pi^2s}\,\frac{p_f}{p_i}\,|\mathcal M|^2.
\end{equation*}

\rtheory{Четырёхвекторы и метрика}

\begin{align*}
    A^\mu&=(A^0,\vec A) && \text{--- \rconcept{контравариантный} вектор},\\
    A_\mu&=(A^0,-\vec A) && \text{--- \rconcept{ковариантный} вектор}.
\end{align*}
Скалярное произведение:
\begin{equation*}
    A\cdot B=A^0B^0-\vec A\vec B
    =A_\mu B^\mu=A^\mu B_\mu
    =g_{\mu\nu}A^\mu B^\nu
    =g^{\mu\nu}A_\mu B_\nu.
\end{equation*}
\rconcept{Метрический тензор}
\begin{equation*}
    g_{\mu\nu}=g^{\mu\nu}=
    \begin{pmatrix}
        1&0&0&0\\
        0&-1&0&0\\
        0&0&-1&0\\
        0&0&0&-1
    \end{pmatrix},
    \qquad
    g_{\mu\nu}g^{\mu\nu}=4.
\end{equation*}

\rtheory{Правила Фейнмана для бесспиновых частиц}

\subsection*{Рассеяние \texorpdfstring{$e^-\mu^-\to e^-\mu^-$}{e-mu- → e-mu-}}

\rassumption{Пока $s=0$}: электрон и мюон считаем бесспиновыми.
Обмен фотоном изображается диаграммой на рисунке.

\begin{rbox}[left=16pt,right=16pt,top=12pt,bottom=10pt]{[]}
    \centering
    \begin{tikzpicture}[>=Latex,every node/.style={text=rtext}]
        \coordinate (V1) at (0,1);
        \coordinate (V2) at (0,-1);
        \draw[feyn fermion] (-2.4,1) -- (V1)
            node[pos=0.45,above] {$p_A$};
        \draw[feyn fermion] (V1) -- (2.4,1)
            node[pos=0.55,above] {$p_C$};
        \draw[feyn fermion] (-2.4,-1) -- (V2)
            node[pos=0.45,below] {$p_B$};
        \draw[feyn fermion] (V2) -- (2.4,-1)
            node[pos=0.55,below] {$p_D$};
        \draw[feyn photon] (V1) -- (V2);
        \draw[->,rtextmuted] (0.35,0.55) -- (0.35,-0.55)
            node[midway,right] {$q$};
        \node[left] at (-0.05,0) {$\gamma$};
        \node[left] at (-2.4,1) {$e^-$};
        \node[right] at (2.4,1) {$e^-$};
        \node[left] at (-2.4,-1) {$\mu^-$};
        \node[right] at (2.4,-1) {$\mu^-$};
        \fill[rtext] (V1) circle(1.6pt);
        \fill[rtext] (V2) circle(1.6pt);
    \end{tikzpicture}
    \captionsetup{font=footnotesize,labelfont=bf,skip=5pt}
    \captionof{figure}{Диаграмма Фейнмана для рассеяния $e^-\mu^-\to e^-\mu^-$
    с обменом фотоном.}
    \label{fig:emu_scattering_sem6}
\end{rbox}

Амплитуда:
\begin{equation*}
    -i\mathcal M=
    \underbrace{ie(p_A+p_C)^\mu}_{\substack{\text{вершинный}\\\text{множитель }e^-}}
    \underbrace{\left(\frac{-ig_{\mu\nu}}{q^2}\right)}_{\substack{\text{пропагатор}\\\gamma}}
    \underbrace{ie(p_B+p_D)^\nu}_{\substack{\text{вершинный}\\\text{множитель }\mu^-}}.
\end{equation*}
Из сохранения 4-импульса в вершине
\begin{equation*}
    p_A=q+p_C
    \quad\Longrightarrow\quad
    q=p_A-p_C.
\end{equation*}

\subsection*{Элементы диаграмм}

\begin{center}
    \renewcommand{\arraystretch}{1.4}
    \begin{tabular}{>{\raggedright\arraybackslash}m{0.32\textwidth}
        >{\centering\arraybackslash}m{0.25\textwidth}m{0.25\textwidth}}
        \hline
        \rconcept{\bfseries Элемент} & \rconcept{\bfseries Диаграмма}
        & \rconcept{\bfseries Множитель}\\
        \hline
        Вершина: частица с $s=0$, фотон с $s=1$ &
        \begin{tikzpicture}[>=Latex,every node/.style={text=rtext},baseline=(V)]
            \coordinate (V) at (0,0);
            \draw[feyn fermion] (-0.9,0.6) -- (V) node[pos=0,left] {$p$};
            \draw[feyn fermion] (V) -- (0.9,0.6) node[pos=1,right] {$p'$};
            \draw[feyn photon] (V) -- (0,-0.8);
            \fill[rtext] (V) circle(1.4pt);
        \end{tikzpicture}
        & $ie\,(p+p')^\mu$\\
        Пропагатор фотона, $s=1$ &
        \begin{tikzpicture}[>=Latex,every node/.style={text=rtext},baseline=(A)]
            \coordinate (A) at (-0.9,0);
            \draw[feyn photon] (A) -- (0.9,0);
            \node[below] at (0,-0.1) {$q$};
        \end{tikzpicture}
        & $\dfrac{-ig^{\mu\nu}}{q^2}$\\
        Вершина античастицы ($e^+$) &
        \begin{tikzpicture}[>=Latex,every node/.style={text=rtext},baseline=(V)]
            \coordinate (V) at (0,0);
            \draw[feyn fermion] (V) -- (-0.9,0.6) node[pos=1,left] {$p$};
            \draw[feyn fermion] (0.9,0.6) -- (V) node[pos=0,right] {$p'$};
            \draw[feyn photon] (V) -- (0,-0.8);
            \fill[rtext] (V) circle(1.4pt);
        \end{tikzpicture}
        & $ie\,(-p-p')^\mu$\\
        Аннигиляция $e^-$ с импульсом $p$ и $e^+$ с импульсом $p'$ &
        \begin{tikzpicture}[>=Latex,every node/.style={text=rtext},baseline=(V)]
            \coordinate (V) at (0,0);
            \draw[feyn fermion] (-0.8,0.65) -- (V) node[pos=0,left] {$p$};
            \draw[feyn fermion] (V) -- (-0.8,-0.65) node[pos=1,left] {$p'$};
            \draw[feyn photon] (V) -- (1.1,0);
            \fill[rtext] (V) circle(1.4pt);
        \end{tikzpicture}
        & $ie\,(p-p')^\mu$\\
        \hline
    \end{tabular}
\end{center}
Стрелка на линии показывает направление движения заряда частицы;
у античастицы она \rnote{направлена против её импульса}.

\task{Двухфотонный распад \texorpdfstring{$\eta$}{η}-мезона}
\label{task:eta_gg_sem6}

\remark
Условие этой задачи в записи не сохранилось; от решения остался только фрагмент.

\solution
Энергия $\eta$-мезона делится между двумя фотонами:
\begin{equation*}
    E_1+E_2=E.
\end{equation*}
Инвариантная масса пары фотонов:
\begin{equation*}
    m_\eta^2=2E_1E_2(1-\cos\theta).
\end{equation*}
При $\cos\theta=0$ отсюда
\begin{equation*}
    m_\eta^2=2E_1E_2=2E_1(E-E_1).
\end{equation*}

\end{document}
